% or named: Long Video Generation Architectural Taxonomy
\section{Long Video Generation Architectural Styles}
\label{sec:taxonomy}
We organize video‐generation research into 6 distinct taxonomies (Fig.~\ref{fig:main-taxonomy}), and for each taxonomy tree we provide a detailed analysis along with principled guidelines for selecting the most suitable architectural paradigm. For example, a category like Keyframes-to-Video focuses on generating the video through keyframes images. While a category like discrete temporal chunks, focuses more on dividing the generation into multiple smaller generations stitched together.

\subsection{Keyframes-to-Video}
\label{sec:arch-keyframe}
Unlike one‐shot text‐to‐video methods, several recent works~\cite{chen2025hunyuanvideoavatar,movieagent2025,zhou2024storydiffusion,zhang2023i2vgenxl,liu2023autostory,ho2022vdm} employ a two‐stage generation pipeline. For example, KeyVID~\cite{wang2025keyvid}, unlike others, uses audio-to-video modality. It first partitions the audio stream into audio latent keyframe segments and then synthesizes video for each segment. Likewise, StoryDiffusion~\cite{zhou2024storydiffusion} decomposes the narrative text into a series of sub‐prompts, generates a keyframe image per sub‐prompt, and uses a motion‐prediction module to interpolate these keyframes into a continuous video sequence. This two‐stage paradigm enables scalable long‐duration video synthesis by producing keyframes, followed by motion interpolation or generation, and stitching the resulting segments into arbitrarily long videos. This methodology ensures the consistency of the whole video semantically. However, the sequential nature of keyframe and video synthesis incurs extra latency compared to end‐to‐end approaches, and requires distinct text‐to‐image and text‐to‐video models when a single model does not natively support both modalities.

%%%%%%%%%%%%%%%%%%%%%%%%%%%%%%%%%%%%
%%%%%%%%%%%%%%%%%%%%%%%%%%%%%%%%%%%%

\subsection{Discrete Temporal Chunks}
\label{sec:arch-chunks}
In this approach, a video is partitioned into disjoint temporal chunks of N frames (e.g., 8, 16, or 25). Each chunk is generated independently and then concatenated to reconstruct the full sequence. By capping memory usage at the chunk level, this scheme significantly reduces peak GPU requirements and naturally supports chunks parallel processing. The primary drawback is the potential for artifacts at chunk boundaries. Several recent studies employ this chunk-based paradigm \cite{teng2025magi1, gao2024ca2vdm, chen2023seine}. For example, MAGI-1 \cite{teng2025magi1} partitions videos into 24-frame segments and holistically denoises each one. The denoised output then conditions the subsequent segment, with four segments processed concurrently. CA2-VDM \cite{gao2024ca2vdm} further demonstrates that such chunked training schedules often require additional epochs to learn the diverse chunked boundaries across segments.

%%%%%%%%%%%%%%%%%%%%%%%%%%%%%%%%%%%%
%%%%%%%%%%%%%%%%%%%%%%%%%%%%%%%%%%%%

\subsection{High Compression}
\label{sec:arch-compression}
Most existing video-generation models demand high-end GPUs at inference. To address this, several recent works~\cite{zhang2025framepack, hacohen2025ltxvideo, lin2024opensoraplan, jin2024pyramidflow} have aggressively reduced model size and parameter count. LTXVideo~\cite{hacohen2025ltxvideo} introduces Video-VAE, a variational autoencoder that compresses spatiotemporal dimensions by 192× into a 128-channel latent tensor without patchification, drastically reducing token count and enabling low-latency inference; however, this level of compression sacrifices fine-grained texture details and subtle motions, and can introduce artifacts in regions of rapid movement. FramePack~\cite{zhang2025framepack} enforces a fixed context length via relative temporal weighting, assigning 50\% to the most recent frame, 25\% to the prior, 12.5\% to the next, etc. This supports efficient extension across arbitrary durations but often yields outputs with limited diversity in background and motion.

%%%%%%%%%%%%%%%%%%%%%%%%%%%%%%%%%%%%
%%%%%%%%%%%%%%%%%%%%%%%%%%%%%%%%%%%%
\subsection{Flattened 3D Space-Time (One-Shot)}
\label{sec:arch-one-shot}
The flattened 3D space–time one-shot approach regards the entire video as a single spatiotemporal tensor and synthesizes it in a single forward pass. By folding the temporal axis into a unified spatial latent and applying 3D convolutional diffusion-transformer blocks, these models jointly capture cross-frame dependencies to produce high-fidelity clips of fixed duration as in ~\cite{ho2022vdm, yu2023magvitv2}. However, this end-to-end formulation imposes heavy GPU requirements, which in turn limit achievable clip length and spatial resolution~\cite{wan2025}. Crucially, it ensures semantic coherence across frames while effectively modeling long-range motion correlations. Most of the video generation models are using one-shot techniques. However, the one-shot models differ in their focus as in WAN2.1 \cite{wan2025} aiming to build a foundational model. On the other side, phantom \cite{liu2025phantom} focuses more on subject-personalization. 
%%%%%%%%%%%%%%%%%%%%%%%%%%%%%%%%%%%%
\subsubsection{Foundational One-Shot}
\label{sec:arch-found}
Foundational one-shot models formalize the essential blueprint for treating the entire video as a unified spatiotemporal tensor and generating it in a single pass. They learn a joint prior over the full video tensor via 3D UNet Diffusion, DiT, or MM-DiT backbones ~\cite{unet2015, dit2023, flowmatching2023, flowmatching2024}, enabling one-shot synthesis of videos. That is beside using a 3D VAE that compresses each clip and converts it into a dense latent grid. Papers like ~\cite{kong2024hunyuanvideo, wan2025, peng2025opensora2, vace2025, yu2023magvitv2, bartal2024lumiere, ho2022vdm} focus on creating Foundational models. For instance, WAN2.1 ~\cite{wan2025} uses a single DiT instead of using MM-DiT; however it complements by using a cross-attention module. While HunyuanVideo ~\cite{kong2024hunyuanvideo}, employs a dual-stream methodology MM-DiT that converts separate text and video streams into a flattened 3D space-time. 
%%%%%%%%%%%%%%%%%%%%%%%%%%%%%%%%%%%%
\subsubsection{Single-Subject Personalization}
\label{sec:arch-single}
Single-subject personalization methods adapt a base one-shot video generator to faithfully reproduce a target individual’s appearance given only a one or a few reference frame as in \cite{liu2025phantom,polyak2024moviegen,gao2023dreamvideo,gu2023animatediff,yang2024cogvideox}. They achieve this by injecting or fine-tuning compact identity modules such as: (i) Embedding adapters inject lightweight projection modules into a frozen generator to encode subject appearance. Phantom~\cite{liu2025phantom} fine-tunes these adapters on a handful of exemplar frames to capture identity style, while MovieGen~\cite{polyak2024moviegen} applies per-subject adapters to preserve facial details across temporal generation. (ii) Textual inversion networks optimize new pseudo-token embeddings that encapsulate personalized identity concepts. DreamVideo~\cite{gao2023dreamvideo} learns these embeddings from static images to steer motion-conditioned synthesis, and CogVideoX~\cite{yang2024cogvideox} extends inversion across video frames for enhanced temporal consistency. (iii) LoRA layers ~\cite{hu2021lora} insert trainable low-rank matrices into attention modules, enabling efficient adaptation of large backbones. AnimateDiff~\cite{gu2023animatediff} leverages LoRA adapters in both U-Net and transformer blocks to personalize motion and appearance with minimal compute, avoiding full-model retraining.

%%%%%%%%%%%%%%%%%%%%%%%
\subsubsection{Multiple-Subject Personalization}
\label{sec:arch-one-multi}
Multi-subject personalization extends one-shot video synthesis to scenes with multiple distinct entities by integrating dedicated modules that encode and fuse each subject’s identity separately as in \cite{chen2025videoalchemist, chen2025skyreelsv2, huang2025conceptmaster, hu2025hunyuancustom}. Multi-subject personalization can be organized into four modular strategies: (i) Cross-Attention Fusion inserts per-entity attention heads into a frozen diffusion backbone so that each subject’s image and text descriptor attend separately. Video Alchemist binds each reference via dedicated cross-attention layers to support open-set multi-entity conditioning without fine-tuning~\cite{chen2025videoalchemist}. (ii) Element Embedding Fusion jointly encodes diverse scene elements into a unified latent space and injects them via fusion modules for coherent multi-entity synthesis. SkyReels-v2 learns an image–text joint embedding that precisely assembles multiple characters and backgrounds in one pass~\cite{chen2025skyreelsv2}. (iii) Decoupled Concept Embeddings learns independent latent vectors for each subject to avoid identity crosstalk, injecting them through diffusion-transformer adapters. ConceptMaster enforces strong per-entity disentanglement by adapting low-rank concept tokens~\cite{huang2025conceptmaster}. (iv) Multi-Modal Adapter Fusion layers modality-specific adapters (text, image, audio) into a unified fusion pipeline to preserve subject consistency across modalities. HunyuanCustom uses LLaVA-based text–image fusion and AudioNet adapters to maintain coherent identities in multi-subject, multi-modal videos~\cite{hu2025hunyuancustom}.

%%%%%%%%%%%%%%%%%%%%%%%%%%%%%%%%%%%%
\subsubsection{Multi-Shot Narrative Planning}
\label{sec:arch-narrative}
Multi-shot narrative planning refers to the generation of long video clips explicitly segmented into discrete shots or scenes, ensuring coherent transitions and consistent visual elements across cuts as in \cite{gao2025seedance, chen2025skyreelsv2, huang2025stepvideo}. (i) End-to-end native planning as in Seedance~\cite{gao2025seedance}. This method integrates shot segmentation directly within its diffusion-transformer backbone via Multishot MM-RoPE and per-shot captions, producing all shots in a single pass with implicit cut boundaries. (ii) Planner-on-top architecture as in StepVideo~\cite{huang2025stepvideo}. This methodology adds a StoryAnchors layer on top of its Step-Video-T2V backbone; an LLM expands the prompt into a script, StoryAnchors predicts key “anchor” frames for each shot, and the base model animates between anchors to yield a coherent multi-shot sequence. (iii) LLM-directed autoregressive planning as in SkyReels-V2~\cite{chen2025skyreelsv2}. This architecture uses a multimodal language model to decompose the user prompt into a sequence of shot instructions, then a diffusion-forcing autoregressive generator renders shots one after another—supporting hard cuts, soft dissolves, and infinite-length film generation.


%%%%%%%%%%%%%%%%%%%%%%%%%%%%%%%%%%%%
%%%%%%%%%%%%%%%%%%%%%%%%%%%%%%%%%%%%

\subsection{Token-Stream Autoregressive-Token}
\label{sec:arch-token}

This method formulates video generation as next-token prediction over a unified text–video token stream. Specifically, a decoder-only transformer with causal attention autoregressively predicts each token conditioned on all previously generated tokens as in \cite{wang2024loong, singer2023videopoet}. VideoPoet~\cite{singer2023videopoet} adopts MagViT-v2~\cite{yu2023magvitv2} to tokenize video clips, whereas Loong \cite{wang2024loong} leverages a causal 3D-CNN encoder followed by vector quantization clustering to produce discrete video tokens. Both then apply a decoder-only transformer to predict the subsequent token in the combined sequence and finally employ super-resolution to recover spatial details. Despite their flexibility, these approaches incur high-frequency detail loss and compression artifacts in fine details. The need to attend over all prior tokens imposes substantial memory and compute overhead, leading to slow inference. Moreover, early tokens are harder to predict and error accumulates across which further degrades long-range temporal coherence and visual fidelity.

%%%%%%%%%%%%%%%%%%%%%%%%%%%%%%%%%%%%%%%
%%%%%%%%%%%%%%%%%%%%%%%%%%%%%%%%%%%%%%%

\subsection{Closed Source Video Generation}
\label{sec:arch-closed}
Proprietary video generation systems have pushed the boundaries of realism and complexity. Kling2.1~\cite{kling2p1}, Runway Gen-3~\cite{runwaygen3}, MiniMax Hailuo~\cite{hailuo2025}, Pika 2.2~\cite{pika1p5}, Sora~\cite{openai2024sora}, and MovieGen~\cite{polyak2024moviegen} have established a substantial performance gap between closed-source and open-source approaches. More recently, Google’s Veo3~\cite{deepmind2025veo3} and ByteDance’s Seedance 1.0~\cite{gao2025seedance} generate high-fidelity videos that adhere to physical laws, support complex multi-subject scenes, capture intricate pose dynamics and cinematic camera movements, and maintain character consistency. 
%These advances bring video generation to the level of photorealistic videos.
